\todo{this whole section}

% Glacial driven sea level change

This report investigates and provides an improved novel method to estimate modern sea level from ice sheet melt, based on sea level physics.

Sea Level (\textbf{SL}), also known as Relative Sea Level, is the difference in elevation between the sea surface and the solid earth surface. A change in this elevation difference is termed Sea Level Change (\textbf{SLC}).

Sea Surface Height (\textbf{SSH}), also known as Absolute Sea Level, is defined as the elevation of the sea surface above a reference ellipsoid. Ocean satellite altimetry measures Sea Surface Height, as it measures the height difference between itself on an orbital ellipsoid and the sea surface.

 However, the quantity of primary interest is SLC, as this is what is directly relevant to coastal communities and policy-makers.




% Importance

Understanding modern SLC is important for blah. Global Mean Sea Level (\textbf{GMSL}) is the areal average height of SL \autocite{hortonMappingSeaLevelChange2018}. GMSL is a particularly important quantity as it decreases the short timescale effects that affect SL, and allows for tracking of global trends. It is genererally accepted that GMSL change is to a large part influenced by ice sheet and glacial mass flux signals \autocite[e.g.][]{slangenEvaluatingModelSimulations2017}.


Smth regional SLC


% Idea

Ice melt driven SLC has an rotationally vairant impact on the sea surface.

This means that as well as being able to calcualte GMSL from SLC, it is possible to calcualte where the ice sheet change occured to drive the SLC.

Can be useful to understand where melt is occuring etc

% traditional method

take a surface average of SSH and apply correction (?) and call it SLC

% Problem

errors assocaited with traditional are poorly constrained and come from two main places.

first is there is the intrinsic

also, other assumptive corrections are needed to account for non-ice driven GMSL changes, such as.

also, investigation into density variations haven't been investigate


the errors cannot be accoutned for with a simple correction as they are depenednt on the ice source location

believed that the errors are sigificant and locationally variable


% Solution

By using a framework presented in \textcite{heathcoteScaleableBayesianFramework2026}, it is possible to explore modern ice driven sea level change using the full account of the physics related to the gravitationally consistent sea level changein a infinite-demiensional setting. This allows for novel quantification of the errors assocaited with traditional methods of GMSL


REF DA has shown that using a Bayseian inversion method with tide gauge data can lead to more

as includes physical aspects of the problem built into the problem

by including the physics of the rebound effect and other noise signals belived to increase accuracy

\subsection{Sea level physics in relation to ice sheet change}\label{sec:slp}

This will act as an introduction to how Sea Level is affected by ice sheet change.

SL defines the sea surface as an equipotential surface:

\begin{equation}
	\Phi(\mathbf{x} + \SL \hat{\mathbf{r}}) = \Phi_0 \label{eq:sea_level_equipotential}
\end{equation}

When this is perturbed by an ice thickness change, $\thicknesschangeice$, three things occur simultaneously: the gravitational potential itself changes, by $\phi$; the solid Earth surface, $\mathbf{x}$, moves radially, by $\mathrm{u}$; the value of the potential at the ocean surface shifts, from $\Phi_0$ to $\Phi_0 + \Phi_g$. This results in the SL changing, from $\SL$ to $\SL + \Delta \SL$, to maintain the equipotential condition. This can be captured within the following equation:

\begin{equation}
	(\Phi + \phi)(\mathbf{x} + \mathbf{u} + (\SL + \slc)\mathbf{\hat{r}}) = \Phi_0 + \Phi_g
\end{equation}

If we apply a first order Taylor expansion to the left-hand side, we get

\begin{equation}
  \Phi(\mathbf{x} + \SL \hat{\mathbf{r}}) + \phi + \nabla \Phi \cdot (\mathbf{u} + \slc \hat{\mathbf{r}}) = \Phi_0 + \Phi_g
\end{equation}


which can be simplfied: using \autoref{eq:sea_level_equipotential} we can cancel these terms from both sides; $\nabla\Phi$ represents the gravitational acceleration vector with its radial magnitude is $g$, therefore $\nabla\Phi \cdot \mathbf{\hat{r}} = g$; term $u \cdot \nabla\Phi$ represents the potential change due to the vertical motion of the solid surface:

\begin{equation}
	\phi + u \cdot \nabla\Phi + g \slc = \Phi_g
\end{equation}

Which we can rearrange to get the change in sea level:

\begin{equation}
	\slc = -\frac{1}{g}(u \cdot \nabla\Phi + \phi) + \frac{\Phi_g}{g} \label{eq:sl-equation}
\end{equation}

When discussing SSH, we need to account for the fact that SSH is the height of the sea surface above a reference ellipsoid, which is not an equipotential surface. This means that the change in SSH, $\sshc$, is related to the change in SL, $\slc$, by:

\begin{equation}
  \sshc = \slc + \frac{1}{g}(\mathbf{u} \cdot \nabla\Phi + \phi) \label{eq:sshc_slc_relation}
\end{equation}

where $\frac{1}{g}(\mathbf{u} \cdot \nabla\Phi)$ accounts for the vertical crustal motion and $\frac{1}{g}\phi$ accounts for the change in gravitational potential. Using \autoref{eq:sl-equation} this results in:

\begin{equation}
  \sshc = \frac{\Phi_g}{g}
\end{equation}



If we have an ice load, $\zeta$, which is due to the change in ice thickness, $\thicknesschangeice$ with density $\densityice$, then the ice load is:

\begin{equation}
	\zeta = \densityice (1-C) \thicknesschange{I}
\end{equation}

Where $1 - C$ accounts for the possibility of floating ice and C is the ocean function, which is 1 if the point is ocean and 0 if the point is land. This means that the ice load only applies to points where there is no floating ice, as floating ice does not contribute to the load on the solid earth.





\todo{fix this}
The sea surface, when ignoring shorter term dynamic effects, lies on an equipotential gravity surface. Ice mass change results in two effects to this equipotential surface, firstly the added water volume shifts the surface from one equipotential level to another, and secondariliy the change in mass results in

This can be captured within the following equation



Ocean function is
\begin{align}
	\mathrm{C}(\mathbf{x})= \begin{cases}
		1, &\quad\rho_w\mathrm{SL}(x)-\rho_i\mathrm{I}(x)>0 \\
		0, &\quad\rho_w\mathrm{SL}(x)-\rho_i\mathrm{I}(x)\le0
			\end{cases}
\end{align}

SLC can then be defined as:

\begin{align}
	\Delta SL = -\frac{\Delta M_\text{ice}}{\rho_w A}
\end{align}

Where:
\begin{align*}
	A &= \int_{\partial M}\mathrm{C}\, dS \\ &= r^2\int_{\mathbb{S}^2}\mathrm{C}\, dS
\end{align*}


For Glacial isostatic adjustment, relative sea level history, ice thickness history, and mantle viscosity \citetitle{royICE7G2018} \autocite{royICE7G2018} model is used see \textcite{royGlacialIsostaticAdjustment2015} and \textcite{al-attarPySLFPPythonSea2025} for more details.
